\section{Molecular Verifier}
\label{app:verifier}

The Molecular Verifier is a distributed reward computation system that
evaluates language model completions across three classes of cheminformatics
tasks: de~novo molecular generation, molecular property prediction, and
retrosynthetic planning.  It is deployed as an asynchronous HTTP service built
on FastAPI~\citep{fastapi} with Ray~\citep{ray} providing the distributed
execution backend.  This appendix describes its internal architecture,
task-routing logic, and the reward functions used for each task class.

%----------------------------------------------------------------------
\subsection{System Architecture}
\label{app:verifier:architecture}

\paragraph{Ray actors.}
At startup, a \textsc{MolecularVerifier} object is instantiated as a Ray
remote actor.  Ray enables the verifier to run concurrently with the language
model inference process on separate CPU and GPU resources, and to manage
independent worker pools for computationally intensive operations such as
molecular docking.

\paragraph{Asynchronous request buffering.}
Reward queries arrive concurrently from many rollout workers.  Rather than
evaluating each query independently—which would prevent GPU utilisation from
reaching an efficient level—a \textsc{RewardBuffer} component accumulates
incoming requests for a configurable window $\Delta t$ (default
\SI{20}{\second}).  After this window has elapsed, all pending queries are
merged into a single batch and forwarded to the Ray actor.  This batching
strategy amortises the fixed overhead of GPU kernel launches and reduces
context-switching costs.  Formally, let $\{q_1, \ldots, q_N\}$ be the set of
queries arriving within $[t, t+\Delta t]$.  Their completions
$c_1,\ldots,c_N$ and associated metadata
$m_1,\ldots,m_N$ are concatenated into a joint batch, scored in a single
forward pass, and then redistributed to the individual futures awaiting each
$q_i$.

%----------------------------------------------------------------------
\subsection{Task Routing}
\label{app:verifier:routing}

Each query is accompanied by a structured metadata object specifying the list
of molecular properties to evaluate, the optimisation objectives, and
numerical target values (when applicable).  The \textsc{MolecularVerifier}
inspects each metadata record and routes the corresponding completion to one
of three sub-verifiers:

\begin{enumerate}
  \item \textbf{GenerationVerifier} — activated when the metadata contains
        property names drawn from the set of supported molecular descriptors
        or docking targets.
  \item \textbf{MolPropVerifier} — activated when the metadata specifies a
        regression or classification objective over a single molecular
        property, requiring the model to predict a numerical value rather than
        generate a novel molecule.
  \item \textbf{ReactionVerifier} — activated when the metadata describes a
        retrosynthetic or forward-reaction task, requiring the model to
        propose a sequence of reaction steps.
\end{enumerate}

Batches may contain queries belonging to different task types; in this case
the three sub-verifiers are called in parallel and their results are
reassembled in the original query order before being returned.

%----------------------------------------------------------------------
\subsection{Answer Parsing}
\label{app:verifier:parsing}

Before any verification logic is applied, the raw model completion is parsed
to isolate the structured answer region.  Three parsing modes are supported:

\begin{itemize}
  \item \textit{answer\_tags}: Content enclosed in
        \texttt{<answer>}\ldots\texttt{</answer>} delimiters is extracted.
        If the extracted span itself contains a
        \texttt{\textbackslash boxed\{\ldots\}} expression, the boxed content
        is further extracted.
  \item \textit{boxed}: The innermost \texttt{\textbackslash boxed\{\ldots\}}
        expression in the raw completion is extracted directly.
  \item \textit{none}: The full completion is used as-is after removing
        special tokens.
\end{itemize}

The same parsing mechanism is shared by all three sub-verifiers, ensuring
consistent answer extraction regardless of task type.

%----------------------------------------------------------------------
\subsection{Generation Verifier}
\label{app:verifier:generation}

\paragraph{SMILES extraction.}
After parsing, the extracted text is tokenised by whitespace and punctuation
delimiters.  Each token is checked against a character-level filter that
retains only strings compatible with the SMILES alphabet
($[\texttt{A-Za-z0-9=\#:+-[]()/\textbackslash @.\%}]^+$)
and additionally requires at least one carbon atom—either an uppercase
\texttt{C} (aliphatic carbon) or more than two lowercase \texttt{c}
characters (indicative of an aromatic ring).  Remaining
candidates are validated with RDKit~\citep{rdkit}: a candidate is accepted if
and only if \texttt{MolFromSmiles} succeeds, the resulting molecule can be
written to a Mol block, and it does not contain a bridged ring system
(defined as two rings sharing more than two atoms).

If no valid SMILES is found, the reward is set to zero.  If more than one
distinct valid SMILES is extracted from a single completion, the reward is
also set to zero, as the model is expected to generate a single unambiguous
molecule.

\paragraph{Property computation.}
For each valid SMILES, the required molecular properties are evaluated using
a \textsc{ScorerRegistry} Ray actor that lazily initialises per-property
evaluators:

\begin{itemize}
  \item \textit{Docking scores} are computed via GPU-accelerated AutoDock-GPU
        \citep{autodock_gpu}.  Ligands are prepared using
        Meeko~\citep{meeko} (3-D conformer generation followed by
        PDBQT conversion).  Multiple docking actor replicas run concurrently,
        with a configurable number of actors per GPU (default 8), to maintain
        high GPU occupancy when the request buffer contains many distinct
        molecules.
  \item \textit{Drug-likeness and physicochemical descriptors} such as QED
        \citep{qed}, synthetic accessibility (SA)~\citep{sa_score},
        and a broad set of RDKit descriptors (molecular weight, topological
        polar surface area, fraction of \textit{sp}$^3$ carbons, etc.) are
        evaluated directly in-process.
  \item \textit{Bioactivity classifiers} (e.g., GSK3$\beta$, JNK3, DRD2
        inhibition probability) are retrieved from the TDC
        library~\citep{tdc}.
\end{itemize}

\paragraph{Reward normalisation.}
Raw property values are rescaled to $[0, 1]$ using property-specific linear
transformations derived from the distribution of the property over a
reference molecular database.  For docking scores, the rescaling is
$\tilde{r} = (\Delta G + 11) / 10$, which maps a docking score of
$\Delta G = -10~\text{kcal/mol}$ to $\tilde{r}=0.1$ and
$\Delta G = -7~\text{kcal/mol}$ to $\tilde{r}=0.4$.  Because AutoDock-GPU
reports binding free energies as negative numbers (more negative $= $ stronger
binding), a good binder receives a \emph{lower} rescaled value $\tilde{r}$.
The \textit{minimize} objective (see below) is therefore applied to docking
targets so that a stronger binder ultimately receives a higher reward
$r = 1 - \tilde{r}$.

\paragraph{Per-property reward functions.}
Given a rescaled property value $v$ and a target $\tau$, the contribution of
each property to the overall reward depends on the declared objective:
\begin{align}
  r_{\text{below}}  &= \mathbf{1}[v \leq \tau],    \\
  r_{\text{above}}  &= \mathbf{1}[v \geq \tau],    \\
  r_{\text{maximize}} &= v,                         \\
  r_{\text{minimize}} &= 1 - v,                     \\
  r_{\text{equal}} &= \mathrm{clip}\!\left(1 - 100\,(v-\tau)^2,\ 0,\ 1\right).
\end{align}

\paragraph{Multi-property aggregation.}
When a query specifies $K$ properties with individual rewards
$r_1,\ldots,r_K$, the final generation reward is their geometric mean:
\begin{equation}
  \label{eq:geo_mean}
  R = \left(\prod_{k=1}^{K} r_k\right)^{1/K}.
\end{equation}
The geometric mean is preferred over the arithmetic mean because it enforces
that a molecule must achieve a satisfactory score on \emph{all} properties
simultaneously: a zero score on any single property collapses the aggregate
reward to zero regardless of performance on the others.

%----------------------------------------------------------------------
\subsection{Molecular Property Verifier}
\label{app:verifier:molprop}

The molecular property verifier is applied when the model is asked to
\textit{predict} the value of a molecular property, rather than generate a
new molecule.  After parsing, the extracted answer string is scanned for
numerical expressions using a set of regular-expression patterns that handle
plain floats, scientific notation (\textit{e.g.}, $1.3 \times 10^{-5}$,
\texttt{1.3e-5}), HTML \texttt{<sup>} markup, LaTeX notation
(\texttt{\textbackslash{}times 10\textasciicircum\{-5\}}), Unicode superscript
digits, and plus-minus uncertainty expressions
(\textit{e.g.}, $-2.1 \pm 0.5 \to -2.1$).  Range expressions of the forms
``$a$ to $b$'' and ``between $a$ and $b$'' are resolved to their midpoint.

If exactly one numerical value can be unambiguously extracted, the reward is
computed according to the declared objective:

\begin{itemize}
  \item \textit{Regression}: The reward is
        $r = \mathrm{clip}(1 - ((v - \tau) / \sigma)^2,\ 0,\ 1)$, where
        $\sigma$ is a per-property normalisation factor derived from the
        dataset variance.  This penalises deviations from the ground-truth
        value $\tau$ quadratically, rewarding close but imperfect predictions
        and giving zero reward beyond a distance of $\sigma$ from the target.
  \item \textit{Classification}: The reward is $r = \mathbf{1}[\hat{y} = y]$,
        where $y \in \{0, 1\}$ is the ground-truth label and $\hat{y}$ is
        extracted from the answer by matching canonical keywords
        (\textit{true}, \textit{yes}, \textit{1}~$\to$~1; \textit{false},
        \textit{no}, \textit{0}~$\to$~0).
\end{itemize}

Ambiguous or unextractable answers receive a reward of zero.

%----------------------------------------------------------------------
\subsection{Reaction Verifier}
\label{app:verifier:reaction}

The reaction verifier evaluates model responses to retrosynthetic planning
and forward-reaction tasks.  It relies on a pre-computed
\textit{reactant–reaction matrix}, a data structure that encodes, for each
catalogued building block, all reaction templates in which it can participate
and the expected product(s).  This matrix is loaded at server startup and
shared across all Ray workers.

\paragraph{Task types.}
The verifier supports several task configurations:

\begin{itemize}
  \item \textit{Final-product prediction}: The model predicts the product of a
        given reaction; the reward is 1 if the predicted SMILES is identical
        (after canonical form normalisation with RDKit) to the ground-truth
        product, and 0 otherwise.
  \item \textit{Reactant identification}: The model identifies one or more
        reactants required to synthesise a given target; correctness is
        assessed against the ground-truth reactant set.
  \item \textit{Full retrosynthetic path}: The model proposes a multi-step
        synthesis route.  Each proposed reaction step is validated against the
        reaction matrix.  The reward has two components:
        \begin{enumerate}
          \item \emph{Validity}: the fraction of reaction steps that are
                chemically valid according to the reaction matrix,
                $r_{\text{valid}} \in [0, 1]$;
          \item \emph{Product similarity}: either the binary indicator that the
                final product exactly matches the target molecule, or the
                Tanimoto coefficient between the Morgan fingerprints of the
                final product and the target
                (configurable via the \texttt{reaction\_reward\_type}
                parameter).
        \end{enumerate}
        The overall reward is $R = r_{\text{valid}} \cdot r_{\text{product}}$.
  \item \textit{Building-block constrained tasks}: Variants of the above in
        which all proposed reactants must belong to a restricted catalogue of
        commercially available building blocks.  Steps using out-of-catalogue
        building blocks are penalised accordingly.
\end{itemize}
